\section{Face-center thresholds on a simplex}\label{sec:simplex-thresholds}

We retain the symmetric $d$-direction model of Proposition~\ref{prop:dgf}:
the initial direction is uniform, a repeat has probability $p$, and each specified
change has probability $r=(1-p)/(d-1)$. Other ways of modeling a tetrahedral
board ($d=3$) give different laws, which we do not treat.
For $2\le k\le d$ and $N\ge k$, let $C_{N,k}(p)$ denote the probability of
one \emph{balanced face bin}: its $k$ positive coordinates differ by at most
one, and its other coordinates vanish. All such bins have the same probability.
Let $V_N(p)=p^{N-1}/d$ be the probability of one vertex bin, and put
\[
u=\frac{1-p}{(d-1)p},\qquad
a_k=\frac12\log(4\pi)-\frac{k+1/2}{k-1}\log k.
\]
The ratio $C_{N,k}/V_N$, viewed as a function of $u$, depends on $k$ and $N$
but not on the ambient number $d$ of directions.

\begin{theorem}[face-center crossings]\label{thm:simplex-crossing}
For every $2\le k\le d$ and $N\ge k$, there is exactly one
$p=p_{N,k}^{(d)}\in(0,1)$ for which $C_{N,k}(p)=V_N(p)$.
The face center is more likely below this persistence and less likely above it.
The sequence $p_{N,k}^{(d)}$ is strictly increasing in $N$, and
\[
p_{k,k}^{(d)}=\frac1{1+(d-1)(k!)^{-1/(k-1)}}.
\]
Writing $L=\log N$ and
$u_{N,k}=(1-p_{N,k}^{(d)})/((d-1)p_{N,k}^{(d)})$, for fixed $k$ we have
\begin{equation}\label{eq:simplex-root-expansion}
Nu_{N,k}=L-\frac12\log L+a_k
+\frac{\frac14\log L-\frac12a_k+\frac18}{L}
+O_k\!\left(\frac{(\log L)^2}{L^2}\right).
\end{equation}
Consequently, $N(1-p_{N,k}^{(d)})$ equals $(d-1)$ times the right-hand side
of~\eqref{eq:simplex-root-expansion}, with the same error order for fixed $d$.
\end{theorem}

In particular, for the full center $k=d$,
\[
N(1-p_{N,d}^{(d)})=(d-1)\log N-\frac{d-1}{2}\log\log N+c_d+o(1),
\qquad c_d=\frac{d-1}{2}\log(4\pi)-\frac{2d+1}{2}\log d.
\]
For $d=2$, $a_2=\tfrac12\log(\pi/8)$, and the expansion agrees with the
binary one in~\cite{companion}. The proof below does not use the binary
case and treats every face dimension at once.

\begin{theorem}[common critical window]\label{thm:simplex-window}
Fix $d$ and put
\[
u_N(t)=\frac{\log N-\frac12\log\log N+t}{N},\qquad
p_N(t)=\frac1{1+(d-1)u_N(t)}.
\]
Uniformly for $t$ in any fixed compact interval, for each $2\le k\le d$,
\begin{equation}\label{eq:simplex-window}
\frac{C_{N,k}(p_N(t))}{V_N(p_N(t))}
\longrightarrow \exp\bigl((k-1)(t-a_k)\bigr).
\end{equation}
Moreover $a_2>a_3>\cdots$. Thus the face centers cross the vertex height
at distinct constant offsets within a common window of persistence width
$O(1/N)$.
\end{theorem}

We first prove the exact assertions, and then prove a uniform local
estimate from which both asymptotic theorems follow.

\begin{proof}[Proof of the exact assertions in Theorem~\ref{thm:simplex-crossing}]
For a positive count vector $n=(n_1,\ldots,n_k)$, let $W_n(j)$ be the number
of words with those counts and exactly $j$ changes. Its ratio to a vertex
probability is
\[
S_n(u)=\sum_jW_n(j)u^j.
\]
To count $W_n(j)$, choose a proper word of $j+1$ run colors, meaning that
successive colors differ. If color $i$ appears $m_i$ times, its run lengths
can be chosen in $\binom{n_i-1}{m_i-1}$ ways. Thus $W_n(j)$ is a sum of
products of these binomial coefficients. Every coefficient is nonnegative;
the first nonzero one is $W_n(k-1)=k!$. Hence $S_n$ is strictly increasing
from zero to infinity on $u\ge0$ and crosses one exactly once. At $p=0$,
cycling through the $k$ colors and truncating after $N$ positions gives a
word with balanced counts and no successive repeat, so its probability is
positive while the vertex probability is zero. At $p=1$ the reverse strict
inequality holds.

Choose the balanced count vectors as $N$ increases by incrementing one of
their minimum coordinates. Each coefficient $W_n(j)$ is then nondecreasing.
The coefficient of $u^k$ is
\[
W_n(k)=\frac{k!(k-1)}2\sum_i(n_i-1)
=\frac{k!(k-1)}2(N-k).
\]
Indeed, exactly one run color appears twice; for each such color there are
$(k+1)!/2-k!=k!(k-1)/2$ proper arrangements of the multiset of run colors.
This coefficient increases strictly with $N$, so the positive root decreases
strictly with $N$. The substitution $p=1/(1+(d-1)u)$ proves the claimed
monotonicity. If $N=k$, every coordinate equals one and
$S_n(u)=k!u^{k-1}$, giving the initial value. Note also that every root
satisfies $u\le(k!)^{-1/(k-1)}<1$, so all crossings lie in the positive
refresh regime used below.
\end{proof}

\begin{lemma}[positive integral and discrete comparison]\label{lem:simplex-integral}
For fixed $k\ge2$, define
\begin{equation}\label{eq:simplex-F}
F_k(x)=\sum_{m_1,\ldots,m_k\ge1}
M!\prod_{i=1}^k\frac{x^{m_i}}{m_i!(m_i-1)!}
=\int_0^\infty e^{-t}\bigl[\sqrt{xt}\,I_1(2\sqrt{xt})\bigr]^k\,dt,
\quad M=\sum_i m_i.
\end{equation}
Let $n$ be a balanced positive count vector of total $N$, put $h=N/k$,
and, for $0<u<1$, put $x=hu/(1-u)$. With $\mathcal D=x\,d/dx$,
\begin{equation}\label{eq:simplex-discrete-comparison}
0\le 1-\frac{S_n(u)}{(1-u)^{N-1}h^{1-k}F_k(x)/x}
\le\frac{\mathcal D^2F_k(x)}{hF_k(x)}
=O_k\!\left(\frac{(1+x)^2}{h}\right).
\end{equation}
In particular the relative error is $O_k((\log N)^2/N)$ when $x=O(\log N)$.
\end{lemma}

\begin{proof}
The equality in~\eqref{eq:simplex-F} follows from the Bessel power series
and Tonelli's theorem, using $M!=\int_0^\infty e^{-t}t^Mdt$.
The integral is finite for every positive $x$: the bound $I_1(y)\le e^y$
for $y\ge0$ reduces its tail to a polynomial times
$\exp(-t+2k\sqrt{xt})$. Thus the positive series and all its termwise
derivatives converge on compact subsets of $x>0$.

Set $\sigma=p-r=p(1-u)$ and $v=r/\sigma=u/(1-u)$.
Proposition~\ref{prop:dpmf}, restricted to the positive coordinates, gives
\begin{equation}\label{eq:simplex-positive-sum}
S_n(u)=(1-u)^{N-1}\sum_{m_i\ge1}M!v^{M-1}
\prod_i\frac{\binom{n_i-1}{m_i-1}}{m_i!}.
\end{equation}
Define
\[
A_{i,m}=\frac{(m-1)!}{h^{m-1}}\binom{n_i-1}{m-1}
=\prod_{j=1}^{m-1}\left(1-\frac{h-n_i+j}{h}\right)_+.
\]
Since $h-n_i>-1$, these are products of factors in $[0,1]$, including beyond
their finite support. The elementary product-deficit bound in
Lemma~\ref{lem:clipped-product} gives
\[
0\le1-\prod_iA_{i,m_i}
\le\frac1h\sum_i\frac{m_i(m_i+1)}2\le\frac{M^2}{h}.
\]
Replacing the binomial factors in~\eqref{eq:simplex-positive-sum} by
$h^{m_i-1}/(m_i-1)!$ gives the denominator in
\eqref{eq:simplex-discrete-comparison}. Summing the deficit bound proves
its first inequality and the displayed explicit upper bound.

To bound the quotient of derivatives, for $x\ge1$ use
$M^2\le3x^2(1+1/x)^M$. Indeed,
$\log(1+1/x)\ge1/(2x)$ and
$\sup_{y\ge0}y^2e^{-y/2}=16/e^2<3$. Positivity of the coefficients gives
\[
\frac{\mathcal D^2F_k(x)}{F_k(x)}
\le3x^2\frac{F_k(x+1)}{F_k(x)}=O_k(x^2),
\]
where the last step follows from Lemma~\ref{lem:simplex-laplace} below.
For $0<x\le1$, the quotient extends continuously to $x=0$ with value
$k^2$, since $F_k(x)=k!x^k+O_k(x^{k+1})$. This proves the bound.
\end{proof}

\begin{lemma}[Laplace expansion]\label{lem:simplex-laplace}
For each fixed $k\ge2$, as $x\to\infty$,
\begin{equation}\label{eq:simplex-laplace}
F_k(x)=\frac{k^{k/2+1}}{(4\pi)^{(k-1)/2}}
x^{(k+1)/2}e^{k^2x}
\left(1-\frac{k-1}{8kx}+O_k(x^{-2})\right).
\end{equation}
\end{lemma}

\begin{proof}
In~\eqref{eq:simplex-F} set $w=\sqrt t$ and $a=k\sqrt x$.
The exponent becomes $-w^2+2k\sqrt x\,w=k^2x-(w-a)^2$.
First restrict to $|w-a|\le\sqrt x/2$. There $2\sqrt x\,w$ is bounded
below by $(2k-1)x$, so the standard Bessel expansion
\cite[Eq.~10.40.1]{dlmf} applies uniformly:
\[
\bigl[\sqrt x\,w I_1(2\sqrt x\,w)\bigr]^k
=\frac{x^{k/4}w^{k/2}e^{2k\sqrt x\,w}}{(4\pi)^{k/2}}
\left(1-\frac{3k}{16\sqrt x\,w}+O_k(x^{-2})\right).
\]
On the complementary set, the elementary exponential bound for $I_1$
dominates the integrand by
$2x^{k/2}w^{k+1}e^{k^2x-(w-a)^2}$.
Gaussian tail estimates therefore show that its integral is at most
$e^{k^2x-cx}$ times a fixed power of $x$, for some $c>0$.
It is exponentially smaller than the main term.

Put $\alpha=k/2+1$. Apart from that negligible tail, the integral is
\[
\frac{2x^{k/4}e^{k^2x}}{(4\pi)^{k/2}}
\int_{|w-a|\le\sqrt x/2} e^{-(w-a)^2}w^\alpha
\left(1-\frac{3k}{16\sqrt x\,w}+O_k(x^{-2})\right)dw.
\]
Taylor-expand $w^\alpha$ about $a$. The odd terms integrate to zero on this
symmetric interval, the quadratic term contributes
$\alpha(\alpha-1)/(4a^2)$ relatively, and the fourth-order remainder is
$O_k(a^{-4})$: the interval stays a fixed positive fraction of $a$ away
from zero and all Gaussian moments are finite. Extending the centered
Gaussian integrals to the real line changes them exponentially little.
For the term involving $w^{\alpha-1}/\sqrt x$, keeping its constant term
suffices, because its quadratic correction has relative order $x^{-2}$.
Since $\int e^{-y^2}dy=\sqrt\pi$ and
$\int y^2e^{-y^2}dy=\sqrt\pi/2$, the total relative correction is
\[
\frac{\alpha(\alpha-1)}{4k^2x}-\frac3{16x}
=-\frac{k-1}{8kx}.
\]
The leading constant simplifies to the one in
\eqref{eq:simplex-laplace}, completing the proof.
\end{proof}

\begin{proof}[Proof of the asymptotic assertions in
Theorems~\ref{thm:simplex-crossing} and~\ref{thm:simplex-window}]
Set $s=k-1$, $z=Nu$, and $D_k=k^{k+1/2}/(4\pi)^{(k-1)/2}$.
The preceding lemmas imply, uniformly on
$c\log N\le z\le C\log N$ for fixed $0<c<C$,
\begin{equation}\label{eq:simplex-uniform-ratio}
S_n(z/N)=D_k\frac{z^{s/2}e^{sz}}{N^s}
\left(1-\frac{s}{8z}+O_{k,c,C}(z^{-2}+z^2/N)\right).
\end{equation}
Indeed, $x=z/(k(1-u))$, and
\[
k^2x+(N-1)\log(1-u)-sz=O_k(z^2/N).
\]
The prefactor changes by $1+O_k(z/N)$ on replacing $x$ by $z/k$;
rounding of the balanced coordinates is covered by
\eqref{eq:simplex-discrete-comparison}.

At $z=(1-\varepsilon)\log N$ the ratio in
\eqref{eq:simplex-uniform-ratio} tends to zero, whereas at
$z=(1+\varepsilon)\log N$ it tends to infinity. The unique root therefore
has $z/\log N\to1$. Taking logarithms at the root and dividing by $s$ gives
\[
z+\tfrac12\log z=L+a_k+\frac1{8z}+O_k(z^{-2}+z^2/N).
\]
For an explicit inversion, put
\[
d_L=-\tfrac12\log L+a_k,\qquad e_L=\tfrac18-\tfrac12d_L,
\qquad z_*=L+d_L+e_L/L.
\]
Taylor expansion gives
$z_*+\tfrac12\log z_*-1/(8z_*)=L+a_k+
O_k((\log L)^2/L^2)$. The derivative of
$z+\tfrac12\log z-1/(8z)$ is at least one for $z>0$.
The preceding root equation, with $z\sim L$, therefore gives
$z-z_*=O_k((\log L)^2/L^2)$, proving
\eqref{eq:simplex-root-expansion}. The conversion to physical persistence
uses $N(1-p)=(d-1)z/(1+(d-1)z/N)=(d-1)z+O_d(z^2/N)$.
Substituting $z=L-\tfrac12\log L+t$ in
\eqref{eq:simplex-uniform-ratio} proves~\eqref{eq:simplex-window}, uniformly
on compact $t$ intervals.

For the ordering of constants, extend
$g(x)=(x+1/2)\log x/(x-1)$ to $x\ge2$.
The sign of $g'(x)$ is the sign of
$H(x)=2x^2-x-1-3x\log x$.
Here $H(2)=5-6\log2>0$, and
$H'(x)=4x-4-3\log x>0$ for $x\ge2$, since its derivative is
$4-3/x>0$ and $H'(2)>0$. Thus $g$ increases strictly, so $a_k$ decreases
strictly.
\end{proof}

\subsection{Thresholds throughout the interior of a face}

The same argument gives a threshold at every point in the interior of a
face, not just at its center. The constants below are evaluated at the
actual fractions $\alpha_i=n_i/N$, since we do not assume that the count
vector converges at any rate.

\begin{theorem}[interior composition thresholds]\label{thm:simplex-interior}
Fix $k\ge2$ and $\varepsilon>0$. Let $n_1+\cdots+n_k=N$, with
$\alpha_i=n_i/N\ge\varepsilon$, and define
\[
A=\sum_{i=1}^k\sqrt{\alpha_i},\qquad g=A^2-1,\qquad s=k-1,
\qquad
D(\alpha)=\frac{A^{k/2+1}(\prod_i\alpha_i)^{-3/4}}{(4\pi)^{s/2}}.
\]
Uniformly for these count vectors and $c\log N\le z\le C\log N$, with
fixed $0<c<C$,
\begin{equation}\label{eq:simplex-interior-ratio}
S_n(z/N)=D(\alpha)\frac{z^{s/2}e^{gz}}{N^s}
\left(1+O_{k,\varepsilon,c,C}(z^{-1}+z^2/N)\right).
\end{equation}
Let $u_n>0$ be the unique root of $S_n(u)=1$. Writing $L=\log N$,
\begin{equation}\label{eq:simplex-interior-root}
Nu_n=\frac{s}{g}L-\frac{s}{2g}\log L
-\frac{\log D(\alpha)+(s/2)\log(s/g)}{g}
+O_{k,\varepsilon}\!\left(\frac{\log L}{L}+\frac{L^2}{N}\right).
\end{equation}
These estimates are uniform on every compact subset of the relative
interior of the face. For a fixed target $\alpha$ and integer counts
$n_i=N\alpha_i+O(1)$, the constants may equivalently be evaluated at that
target, without changing the error order.
\end{theorem}

\begin{proof}
The proper-run polynomial used above has nonnegative coefficients and
first coefficient $k!u^{k-1}$ for every positive count vector, which gives
the unique positive root. To obtain the approximation, put
$v=u/(1-u)$ and $y=Nv$, and introduce
\[
G_{\alpha}(y)=\int_0^\infty e^{-t}
\prod_i\left[\sqrt{\alpha_i y t}\,
I_1(2\sqrt{\alpha_i y t})\right]dt.
\]
Replacing each binomial coefficient in~\eqref{eq:simplex-positive-sum}
by $n_i^{m_i-1}/(m_i-1)!$ gives
\[
S_n(u)=\frac{(1-u)^{N-1}}{v\prod_i n_i}\,G_{\alpha}(y)
\left(1+O_{k,\varepsilon}((1+y)^2/N)\right).
\tag{*}
\]
To see that the error is uniform, note that the falling-factorial deficit
is bounded by
$\sum_i m_i(m_i-1)/(2n_i)\le M^2/(2\varepsilon N)$.
The positive-series argument in Lemma~\ref{lem:simplex-integral} bounds
the weighted second moment of $M$ by a constant times $(1+y)^2$.
Uniformity follows from the uniform saddle expansion below, which bounds
$G_{\alpha}(y+1)/G_{\alpha}(y)$ on the compact set
$\{\alpha_i\ge\varepsilon,\ \sum_i\alpha_i=1\}$; for bounded $y$,
continuity and the positive leading coefficient give the same conclusion.

Set $w=\sqrt t$. The exponential factor of the integrand is
\[
\exp(-w^2+2A\sqrt y\,w)
=\exp(A^2y-(w-A\sqrt y)^2).
\]
Every $\alpha_i$ stays bounded away from zero, so the Bessel expansions,
Gaussian tail estimates, and Taylor remainders in
Lemma~\ref{lem:simplex-laplace} are uniform here. Keeping their leading
terms yields
\[
G_{\alpha}(y)=
\frac{A^{k/2+1}(\prod_i\alpha_i)^{1/4}}{(4\pi)^{s/2}}
y^{(k+1)/2}e^{A^2y}\left(1+O_{k,\varepsilon}(y^{-1})\right).
\]
Substitute this into (*) and use
$v\prod_i n_i=yN^{k-1}\prod_i\alpha_i$.
For $u=z/N$ in the stated range, $y=z+O(z^2/N)$ and
$(N-1)\log(1-u)=-z+O(z^2/N)$. This proves
\eqref{eq:simplex-interior-ratio} with the specified constant.

On the compact set of compositions, both $g$ and $D(\alpha)$ have positive
lower and finite upper bounds. Testing the ratio at suitable constant
multiples of $\log N$ brackets the unique root uniformly. Taking logarithms
there gives
\[
gz+\frac{s}{2}\log z=sL-\log D(\alpha)
+O_{k,\varepsilon}(z^{-1}+z^2/N).
\]
First $z=(s/g)L+O_{k,\varepsilon}(\log L)$; one substitution then gives
\eqref{eq:simplex-interior-root}. The derivative of its left-hand side is
at least $g$, bounded uniformly away from zero, so the stated error is
preserved on inversion. Finally, the functions of $\alpha$ appearing in
the formula have bounded derivatives on a slightly smaller compact set;
changing $\alpha$ by $O(1/N)$ changes the expression by $O(L/N)$, which
is absorbed by its error term.
\end{proof}

At a balanced composition, $A=\sqrt k$, $g=k-1$, and
$D(\alpha)=D_k$, recovering the preceding center formula through its
constant term. For $k=2$, $g=2\sqrt{\alpha_1\alpha_2}$, recovering the
binary bulk scale. Cauchy--Schwarz gives $g\le k-1$, with equality only at
the balanced composition, so the leading threshold coefficient $s/g$
is minimized exactly at the center of each face.

\subsection{Balanced bins and global modes}

\begin{theorem}[balancing within a support]\label{thm:simplex-balance}
For $1/d\le p<1$, among bins with a prescribed set of positive coordinates
and total $N$, probability is maximized precisely by the balanced count
vectors. More precisely, transferring one unit from a coordinate $b$ to
another coordinate $a$, where $b\ge a+2$ and $a\ge1$, strictly increases
the bin probability.
\end{theorem}

\begin{proof}
At $p=1/d$ this follows from the ordinary multinomial formula. For
$1/d<p<1$, put $\sigma=p-r>0$ and $v=r/\sigma>0$.
The positive refresh sum has the integral form
\begin{equation}\label{eq:simplex-balance-integral}
P_N(n)=\frac{\sigma^{N-1}}{dv}\int_0^\infty e^{-t}
\prod_i B_{n_i}(vt)\,dt,\qquad
B_n(y)=\sum_{m=1}^n\binom{n-1}{m-1}\frac{y^m}{m!}.
\end{equation}
For each $y>0$, the sequence $B_n(y)$, $n\ge1$, is strictly log-concave.
This is a case of the fact that the binomial transform preserves
log-concavity, and we give a direct proof. Write
\[
\frac{B_n(y)}y=b_{n-1},\qquad
b_j=\sum_{i=0}^j\binom ji a_i,\quad
a_i=\frac{y^i}{(i+1)!},\qquad
c_j=\sum_{i=0}^j\binom ji a_{i+1}.
\]
Pascal's identity gives $b_{j+1}=b_j+c_j$.
The ratio $c_j/b_j$ is the expectation of the strictly decreasing sequence
$a_{i+1}/a_i=y/(i+2)$ under weights proportional to $\binom ji a_i$.
When $j$ increases to $j+1$, the likelihood ratio of the new and old
unnormalized weights on the old support is
$(j+1)/(j+1-i)$, increasing in $i$; additionally a strictly positive weight
appears at $i=j+1$. Thus these weights shift strictly toward larger
indices, so the expectation of the strictly decreasing sequence decreases
strictly. To check this last step, cross-multiply the two expectations and
pair the terms with indices $i<\ell$. Hence
$b_{j+1}/b_j=1+c_j/b_j$ decreases strictly with $j$, proving the asserted
strict log-concavity.

It follows that
$B_{a+1}(y)B_{b-1}(y)>B_a(y)B_b(y)$ whenever $b\ge a+2$.
Apply this inequality pointwise for $y=vt>0$ in
\eqref{eq:simplex-balance-integral} and integrate. Repeated transfers give
the conclusion.
\end{proof}

\begin{corollary}[global modes in the critical window]\label{cor:simplex-modes}
Fix $d$. In the window of Theorem~\ref{thm:simplex-window}, if $t<a_d$,
then for all sufficiently large $N$ the only global maxima are the $d$
vertices. If $t>a_d$, then the only global maxima are the fully supported
balanced bins. These assertions are uniform on compact sets of $t$
separated from $a_d$. At the exact full-center/vertex crossing, for all
sufficiently large $N$ the global maxima are exactly the vertices and the
fully supported balanced bins, with both classes present.
\end{corollary}

\begin{proof}
Theorem~\ref{thm:simplex-balance} reduces the possible modes to balanced
face bins of support sizes $1,\ldots,d$. For $2\le k\le d$ their limiting log ratios to a
vertex are $(k-1)(t-a_k)$. If $t<a_d$, these are all negative. If $t>a_d$,
the full-center ratio exceeds the vertex ratio, and for $2\le k<d$ its limiting log ratio exceeds the $k$-face
one by
\[
(d-k)(t-a_d)+(k-1)(a_k-a_d)>0.
\]
Uniformity in Theorem~\ref{thm:simplex-window} gives both assertions.
At the exact full-center crossing, its window parameter tends to $a_d$
by Theorem~\ref{thm:simplex-crossing}; every proper face-center ratio then
has a limit strictly less than one. The same reduction proves the final
claim.
\end{proof}

\begin{example}[an intermediate face can be a finite-size mode]\label{ex:simplex-edge-mode}
Corollary~\ref{cor:simplex-modes} does not hold for every $N$.
Take $d=3$, $N=4$, and $p=10/17$, so $u=7/20$.
For an edge center and a full center respectively,
\[
\frac{P_4(2,2,0)}{V_4}=2u+2u^2+2u^3=\frac{4123}{4000}>1,
\qquad
\frac{P_4(2,1,1)}{V_4}=6u^2+6u^3=\frac{3969}{4000}<1.
\]
By Theorem~\ref{thm:simplex-balance}, the three edge centers are the
only global modes. So at finite $N$ the modes can lie on an intermediate
face.
\end{example}
